% !TeX program = xelatex
% ============================================================
% main.tex — Probearbeit "Bier selber brauen"
% Aufbau gemäss Handbuch Projekte KSWO 2025, Kap. 7.1.2
% (naturwissenschaftliche Arbeit).
%
% Kompilieren: xelatex main.tex && biber main && xelatex main.tex && xelatex main.tex
% (oder: latexmk main.tex, siehe latexmkrc)
%
% Stilhinweis (Handbuch Kap. 7.3): Zahlen bis zwölf sowie alle
% Zehnerzahlen im Lauftext ausschreiben ("drei Wochen", nicht "3 Wochen").
% Das ist eine Schreibregel und wird hier nicht automatisiert.
% ============================================================

% ============================================================
% preamble.tex
% Formatierung gemäss Handbuch Projekte KSWO 2025, Kap. 7.3 (Tab. 14)
% und Kap. 7.5 (Bibliografieren und Zitieren).
% Jede Vorgabe ist mit der entsprechenden Handbuchstelle kommentiert.
% ============================================================

\documentclass[11pt,a4paper]{scrreprt}

% ---- Papierformat / Seitenränder (Tab. 14) --------------------
% A4; oben 3cm, unten 2.8cm, links 3cm, rechts 2cm;
% Kopfzeile 1.2cm, Fusszeile 1.8cm.
\usepackage[
  a4paper,
  top=3cm, bottom=2.8cm, left=3cm, right=2cm,
  headheight=1.2cm, footskip=1.8cm,
]{geometry}

% ---- Schrift (Tab. 14: "gut lesbare Schrift, z.B. Calibri oder Arial") --
% Carlito ist metrisch kompatibel zu Calibri und liegt lokal bereits vor.
% Alternative für Arial-Look: Liberation Sans (Zeile einfach auskommentieren).
\usepackage{fontspec}
\setmainfont{Carlito}
% \setmainfont{Liberation Sans}

% ---- Sprache, Silbentrennung, Anführungszeichen ----------------
% babel funktioniert auch mit XeLaTeX/fontspec und stellt die
% Silbentrennungsmuster sowie die Sprachstrings für biblatex/caption.
\usepackage[ngerman]{babel}
\usepackage{csquotes}
\setquotestyle{swiss}  % «...» wie im Handbuch selbst verwendet (z.B. S. 7)

% ---- Zeilenabstand (Tab. 14: 1.5 Text / 1.0 Fussnoten & Zitate) --------
\usepackage{setspace}
\onehalfspacing

% Fussnoten: 9pt (Tab. 14). Nur die Schriftgrösse global umdefinieren --
% \footnotesize wird von vielen Paketen intern zur Breitenmessung
% verwendet (z.B. pgfplotstable), \singlespacing dort mit einzubacken
% führte zu einer Endlosschleife ("TeX capacity exceeded"). Der einfache
% Zeilenabstand wird stattdessen gezielt nur auf den eigentlichen
% Fussnotentext angewendet.
\makeatletter
\renewcommand\footnotesize{%
  \@setfontsize\footnotesize{9pt}{10.8pt}%
}
\let\kswo@makefntext\@makefntext
\renewcommand\@makefntext[1]{%
  \begingroup\singlespacing\kswo@makefntext{#1}\endgroup}
\makeatother

% Umgebung für längere (eingerückte) Zitate: kleinerer Zeilenabstand/Schrift
% optional, siehe Handbuch S. 98, Regel f)
\newenvironment{zitat}
  {\begin{quote}\singlespacing\small}
  {\end{quote}}

% ---- Gegen verwaiste Zeilen (Tab. 14 / Checkliste S. 103) --------------
\clubpenalty=10000
\widowpenalty=10000
\displaywidowpenalty=10000

% ---- Kopf-/Fusszeile: Kapitelname oben, Seitenzahl unten (Tab. 14) -----
% scrlayer-scrpage statt fancyhdr, da KOMA-Script-kompatibel (kein
% "fancyhdr vs. KOMA"-Konflikt).
\usepackage[automark]{scrlayer-scrpage}
\clearpairofpagestyles
\ohead{\headmark}
\ofoot[\pagemark]{\pagemark}
\setkomafont{pageheadfoot}{\normalfont\small}

% ---- Abbildungen/Tabellen: fortlaufende Nummerierung übers ganze -------
% Dokument statt pro Kapitel (Handbuch S. 88, "Nummerierungen und Verweise")
\makeatletter
\@removefromreset{figure}{chapter}
\@removefromreset{table}{chapter}
\makeatother
\renewcommand{\thefigure}{\arabic{figure}}
\renewcommand{\thetable}{\arabic{table}}

% Kurzform "Abb."/"Tab." statt "Abbildung"/"Tabelle" (wie im Handbuch
% selbst verwendet, z.B. Abb. 33/34, Tab. 14/15/16)
\usepackage{caption}
% labelformat=simple erzwingen: KOMA-Script hängt in seinem eigenen
% Default sonst zusätzlich einen Punkt an die Nummer an ("Tab. 1."),
% was zusammen mit labelsep=colon zu "Tab. 1.:" führt statt "Tab. 1:"
% wie im Handbuch (z.B. S. 9, "Tab. 1:").
\captionsetup{labelformat=simple, labelsep=colon, font=small, labelfont=bf}
\addto\captionsngerman{\renewcommand{\figurename}{Abb.}}
\addto\captionsngerman{\renewcommand{\tablename}{Tab.}}

% ---- Grafiken, Tabellen, Einheiten (für Messwerte beim Brauen) ---------
\usepackage{graphicx}
\graphicspath{{bilder/}}
\usepackage{booktabs}
\usepackage{siunitx}
\usepackage[table]{xcolor} % für Hintergrundfarben in Tabellen

% ---- Tabellenfarben (Vorlage tabelle-beispiel_2.html) ------------------
% Kopfzeile dunkelgelb/orange mit schwarzer Schrift, Datenzeilen hell
% beige hinterlegt, Rahmenlinien weiss (wie im HTML-Beispiel).
\definecolor{bierheaderbg}{HTML}{E1A339}
\definecolor{bierrowbg}{HTML}{ECC683}

% Tabellen und Diagramme direkt aus CSV-Messdaten erzeugen (Daten bleiben
% in der .csv, werden beim Kompilieren eingelesen statt in den
% .tex-Quellcode kopiert)
\usepackage{pgfplotstable}
\pgfplotstableset{col sep=semicolon, string type}
\usepackage{xstring} % zum Umformatieren von Datumswerten aus der CSV
\usepackage{pgfplots}
\pgfplotsset{compat=1.18}

% Datumswerte im ISO-Format "YYYY-MM-DD HH:MM" (so stehen sie in der
% CSV) werden beim Einlesen auf "DD.MM.YYYY HH:MM" umformatiert (feste
% Zeichenpositionen, passend zum festen CSV-Format). Anwendung:
% columns/<Spalte>/.style={ddmmyyyy from iso}
% Hinweis: keine "@" in den Hilfsmakronamen verwenden -- dieser Block
% liegt ausserhalb von \makeatletter/\makeatother, "@" hätte dort
% Kategoriecode 12 (kein Buchstabe) und würde die Makrodefinitionen
% brechen.
\pgfplotstableset{
  ddmmyyyy from iso/.style={
    /pgfplots/table/preproc cell content/.append code={
      \pgfkeysgetvalue{/pgfplots/table/@cell content}\kswoRawDT
      \StrLeft{\kswoRawDT}{4}[\kswoIsoY]
      \StrMid{\kswoRawDT}{6}{7}[\kswoIsoM]
      \StrMid{\kswoRawDT}{9}{10}[\kswoIsoD]
      \StrMid{\kswoRawDT}{12}{16}[\kswoIsoT]
      \edef\kswoNewDT{\kswoIsoD.\kswoIsoM.\kswoIsoY\space\kswoIsoT}
      \pgfkeyslet{/pgfplots/table/@cell content}\kswoNewDT
    }
  }
}

% ---- Literaturverzeichnis: Kurzbelege in Klammern (Handbuch Kap. 7.5) --
% "Meier 1993, S. 23" bzw. "Meier; Müller; Maurer 2003, S. 13-15" bzw.
% "Meier u.a. 2003, S. 13-15" bei drei und mehr Autoren.
% Hinweis: Die Handbuch-Regeln sind kleinteiliger, als es jeder
% Standardstil abbildet (z.B. exaktes Datumsformat bei Internetquellen,
% Sonderfälle bei Sekundärzitaten). Bei Bedarf hier feinjustieren.
\usepackage[
  backend=biber,
  style=authoryear,
  natbib=false,
  maxcitenames=3,
  mincitenames=1,
  maxbibnames=99,
  uniquename=false,
  uniquelist=false,
  giveninits=false,
  dashed=false,
  date=year,
  urldate=terse,
  datezeros=true,
  sorting=nyt,
]{biblatex}
\addbibresource{quellen.bib}

% Kein Komma zwischen Autor und Jahr: "Meier 1993" statt "Meier, 1993"
\renewcommand*{\nameyeardelim}{\space}
% Semikolon statt "und" zwischen mehreren Autoren (Handbuch S. 95)
\renewcommand*{\multinamedelim}{\addsemicolon\space}
\renewcommand*{\finalnamedelim}{\addsemicolon\space}
% "u.a." statt "et al." (Handbuch S. 100 nennt beide Varianten, "u.a."
% wird hier als Standard gewählt)
\DefineBibliographyStrings{ngerman}{andothers = {u.\,a.}, urlseen = {Abruf}, edition = {Auflage}}

% Titel im Quellenverzeichnis nicht kursiv (Handbuch-Beispiele S. 94-97
% verwenden durchgehend einfachen Text ohne Auszeichnung)
\DeclareFieldFormat[book,collection,online,report,thesis,article,incollection]{title}{#1}
\DeclareFieldFormat{journaltitle}{#1}
\DeclareFieldFormat{booktitle}{#1}
% Bare URL ohne automatisches "URL:"-Präfix (Handbuch-Beispiele S. 96
% zeigen die nackte Adresse)
\DeclareFieldFormat{url}{\url{#1}}

% Tageszahl/Monat des Abrufdatums zweistellig ausgeben (z.B. "06" statt
% "6"), da \thefield{urlmonth} nicht automatisch nullgepolstert wird.
\newcommand*{\kswopad}[1]{\ifnum#1<10 0#1\else #1\fi}

% ---- Eigene Bibliografie-Vorlagen für das Quellenverzeichnis -----------
% Standard-authoryear-Stile setzen "(Jahr)" direkt nach dem Namen; das
% Handbuch will "Nachname, Vorname: Titel. ... Ort Jahr." (S. 94 ff.).
% Hier für die vier häufigsten Fälle nachgebaut: Monografie (@book),
% Internetquelle (@online), Sammelbandaufsatz (@incollection) und
% Zeitschriftenaufsatz (@article). Für Zeitungsartikel, Interviews,
% AV-Medien etc. (Handbuch S. 95-97) nach demselben Muster ergänzen.
\makeatletter

% Monografien (Handbuch S. 94):
% "Nachname, Vorname: Titel. Untertitel. Bd. (Nr.). Auflage. Ort Jahr."
% \printfield{edition} hängt bei numerischer Edition bereits selbst
% ". <Auflagenstring>" an -- hier NICHT nochmals manuell ergänzen.
\DeclareBibliographyDriver{book}{%
  \printnames{author}%
  \setunit{\addcolon\space}%
  \printfield{title}%
  \newunit
  \printfield{subtitle}%
  \newunit
  \printfield{edition}%
  \newunit
  \printlist{location}%
  \setunit*{\addspace}%
  \printfield{year}%
  \newunit\finentry}

% Internetquellen (Handbuch S. 96):
% "Nachname, Vorname (sofern vorhanden): Titel. URL (Abruf TT.MM.JJJJ)."
% Achtung: "author" ist ein Namensfeld -- mit \ifnameundef prüfen, nicht
% mit \iffieldundef (das erkennt Namensfelder fälschlich als leer).
\DeclareBibliographyDriver{online}{%
  \ifnameundef{author}
    {\printfield{title}}
    {\printnames{author}\setunit{\addcolon\space}\printfield{title}}%
  \newunit
  \printfield{url}%
  \iffieldundef{urlyear}{}{%
    \space\printtext[parens]{\bibstring{urlseen}~%
      \kswopad{\thefield{urlday}}.\kswopad{\thefield{urlmonth}}.\thefield{urlyear}}}%
  \newunit\finentry}

% Aufsätze in Sammelbänden (Handbuch S. 95):
% "Nachname, Vorname: Aufsatztitel. In: Nachname Hrsg., Vorname (Hrsg.):
%  Buchtitel. Untertitel. Bd. (Nr.). Ort Jahr, Seiten."
% BibTeX-Typ @incollection, "booktitle"/"booksubtitle" = Titel des
% Sammelbandes, "editor" = Herausgeber/in, "volume" = Bandnummer.
% Hinweis: \printfield{volume} hängt selbst bereits "Bd. " vor die Zahl
% (Standardverhalten von biblatex) -- hier nicht nochmals manuell ergänzen.
\DeclareBibliographyDriver{incollection}{%
  \printnames{author}%
  \setunit{\addcolon\space}%
  \printfield{title}%
  \newunit
  \printtext{In\addcolon\space}%
  \printnames{editor}\addspace\printtext{(Hrsg.):}\space
  \printfield{booktitle}%
  \newunit
  \printfield{booksubtitle}%
  \newunit
  \iffieldundef{volume}{}{\printfield{volume}\newunit}%
  \printlist{location}%
  \setunit*{\addspace}%
  \printfield{year}%
  \setunit{\addcomma\space}%
  \printfield{pages}%
  \newunit\finentry}

% Aufsätze in Zeitschriften (Handbuch S. 95):
% "Nachname, Vorname: Aufsatztitel. In: Zeitschriftentitel mit
%  Bandnummer, Seiten."
% BibTeX-Typ @article. Handbuch-Beispiel "Praxis Geschichte 2/99" ist
% Heft/Jahr in einem Feld -- dafür das "number"-Feld wörtlich so befüllen
% (z.B. number = {2/99}), statt volume/number/year einzeln zu pflegen.
\DeclareBibliographyDriver{article}{%
  \printnames{author}%
  \setunit{\addcolon\space}%
  \printfield{title}%
  \newunit
  \printtext{In\addcolon\space}%
  \printfield{journaltitle}%
  \iffieldundef{number}{}{\setunit{\space}\printfield{number}}%
  \setunit{\addcomma\space}%
  \printfield{pages}%
  \newunit\finentry}

\makeatother

% Kurzbefehl für Paraphrasen mit vorangestelltem "Vgl." (Handbuch S. 97/99)
% Verwendung im Text: \vgl{quellenkey} oder mit Seitenzahl \vgl[S.~45]{quellenkey}
\newrobustcmd*{\vgl}[2][]{\parencite[vgl.][#1]{#2}}

% ---- Verlinkung (praktisch, keine Handbuchvorgabe) ---------------------
\usepackage[hidelinks]{hyperref}

% Hinweis: preamble.tex enthält absichtlich kein \begin{document} --
% es wird per \input in main.tex eingebunden, das \begin{document} steht
% dort.


\begin{document}

% ---- Titelblatt, Inhaltsverzeichnis, Vorwort ----------------------
% Werden bei der Paginierung mitgezählt, zeigen aber keine Seitenzahl
% (Handbuch Tab. 14, "Paginierung": "beginnt nach dem Inhaltsverzeichnis").
\pagenumbering{arabic}
\pagestyle{empty}

% ============================================================
% titelblatt.tex
% Obligatorische Elemente (Handbuch S. 78, Tab. 12):
% Titel (und evtl. Untertitel), Typ der Arbeit, Name der Autorinnen/
% Autoren, Schule und Abteilung, Datum der Abgabe, Referentin/Referent.
% ============================================================

\begin{titlepage}
\centering

\vspace*{2cm}

{\LARGE\bfseries Bier selber brauen\par}
\vspace{0.5em}
{\Large Vom Rohstoff zur Flasche -- Planung, Durchführung und
  Auswertung eines eigenen Brauprozesses\par}

\vspace{3cm}

{\large Probearbeit im Rahmen des Projektunterrichts\par}

\vspace{3cm}

{\large
  Elias König \\
  Aron Pasold \\
}

\vspace{2cm}

{\large
  Kantonsschule Wohlen \\
  G2028B \\
}

\vspace{2cm}

{\large
  % \today setzt automatisch das heutige Kompilierdatum -- durch ein
  % festes Datum ersetzen, sobald das effektive Abgabedatum feststeht.
  Abgabedatum: \today \\
  Betreuende Lehrperson: Sandro Soom \\
}

\vfill

\end{titlepage}


\tableofcontents
\thispagestyle{empty}

% Vorwort ist optional (Handbuch S. 79). Falls nicht gebraucht: Zeile
% auskommentieren und kapitel/00-vorwort.tex löschen.
% ============================================================
% Vorwort (optional, Handbuch S. 79)
% Bei natur-/geisteswissenschaftlichen Arbeiten eher unüblich -- nur
% verwenden, wenn ihr z.B. persönlich erklären wollt, wie ihr zum Thema
% gekommen seid, und wem ihr danken möchtet.
%
% Falls nicht gebraucht: diese Datei leeren (nicht löschen, sonst bricht
% \include in main.tex ab) oder in main.tex die entsprechende Zeile
% auskommentieren und die Datei entfernen.
% ============================================================

\chapter*{Vorwort}
\addcontentsline{toc}{chapter}{Vorwort}
% \chapter* erzwingt sonst die "plain"-Kapitelseite mit sichtbarer
% Seitenzahl -- hier explizit unterdrücken (Handbuch Tab. 14:
% Titelblatt/Inhaltsverzeichnis/Vorwort ohne Seitenzahl).
\thispagestyle{empty}



% ---- Ab hier sichtbare Seitenzahlen --------------------------------
\pagestyle{scrheadings}

% ============================================================
% Abstract (Handbuch S. 83)
% Max. eine A4-Seite, allgemeinverständlich. Fragestellung, Methode,
% Ergebnisse und Diskussion so kurz und prägnant wie möglich zusammen-
% fassen. Keine Literaturhinweise, Quellenangaben oder unübliche
% Abkürzungen.
% ============================================================

\chapter*{Abstract}
\addcontentsline{toc}{chapter}{Abstract}


% ============================================================
% Einleitung -- trägt die Kapitelnummer 1 (Handbuch S. 83)
%
% Inhalt gemäss Handbuch:
% - Untersuchungsgegenstand, Problemstellung, Theorie und Ziel der
%   Arbeit definieren
% - Untersuchungsfeld klar eingrenzen, Ausgangslage darstellen
%   (grösserer Zusammenhang, vorhandene Literatur, bestehende Versuche)
% - Am Schluss der Einleitung: Hypothesen, die sich aus dem Theorieteil
%   ableiten lassen
%
% Für dieses Thema z.B.: Warum Bier selber brauen? Welche Fragestellung
% wird untersucht (z.B. Einfluss von Gärtemperatur/Hefesorte auf
% Geschmack/Alkoholgehalt)? Welche Hypothese wird geprüft?
% ============================================================

\chapter{Einleitung}
\label{kap:einleitung}

% Beispieltext -- als Ausgangspunkt gedacht, noch zu ergänzen und mit
% eigenen Formulierungen zu ersetzen.

Bier zählt zu den ältesten industriell hergestellten Lebensmitteln der
Menschheit und wird bis heute nach einem im Kern unveränderten Prinzip
gebraut: Aus Wasser, Malz, Hopfen und Hefe entsteht durch einen
kontrollierten biochemischen Prozess ein alkoholhaltiges Getränk
\parencite{bruecklmeier2022}. Obwohl der Brauprozess auf den ersten
Blick einfach erscheint, hängt das Ergebnis von einer Vielzahl
kontrollierbarer Grössen ab, etwa der Maischetemperatur, der
Hopfenmenge oder der verwendeten Hefesorte.

In der vorliegenden Arbeit brauen wir in einer Gruppe von zwei Personen
ein eigenes Bier und dokumentieren den gesamten Prozess von der
Zutatenbestellung über den Brauvorgang bis zur Abfüllung in Flaschen.
Im Zentrum steht die Frage, wie sich die Gärtemperatur auf den
Vergärungsgrad und den Geschmack des fertigen Biers auswirkt. Wir
vermuten, dass eine höhere Gärtemperatur zu einem höheren
Vergärungsgrad, aber auch zu stärker ausgeprägten Fruchtnoten führt
\vgl[S.~45]{bruecklmeier2022}.


% ============================================================
% Material und Methoden (Handbuch S. 84)
%
% Zuerst Material, dann Methoden beschreiben:
% - Sind die Schritte wiederholbar? Welche Geräte wurden eingesetzt?
% - Woher stammen die verwendeten Zutaten/Chemikalien?
% - Welchen Umfang haben die untersuchten Stichproben?
%
% Für dieses Thema z.B.:
% - Zutatenliste und Bestellung (Malz, Hopfen, Hefe, Wasser,
%   Bezugsquellen/Lieferanten)
% - Braugeräte und -ausrüstung
% - Rezept und Brauprotokoll (Maischen, Läutern, Hopfenkochen, Kühlen,
%   Gären, Abfüllen), inkl. Mengen- und Temperaturangaben
% - Messmethoden (z.B. Refraktometer/Spindel für Stammwürze, pH-Meter)
% ============================================================

\chapter{Material und Methoden}
\label{kap:material-methoden}


% ============================================================
% Resultate (Handbuch S. 84)
%
% Nur Ergebnisse klar darstellen, hier noch NICHT diskutieren.
% Jede Abbildung/Tabelle braucht eine Legende und wird im Fliesstext
% kommentiert (Handbuch S. 88 f.).
%
% Für dieses Thema z.B.:
% - Messwerte aus dem Brauprozess (Stammwürze/°Plato, pH-Wert,
%   Vergärungsgrad, errechneter Alkoholgehalt)
% - Zeitlicher Verlauf der Gärung (z.B. als Tabelle/Abbildung)
% - Verkostungsergebnisse (falls eine Umfrage/Degustation gemacht wurde)
%
% Tabelle und Diagramm unten werden beim Kompilieren direkt aus
% daten/gaerung.csv erzeugt (pgfplotstable/pgfplots) -- die Messwerte
% stehen NUR in der CSV-Datei, nicht im .tex-Quellcode. Bei neuen
% Messungen einfach die CSV aktualisieren.
%
% Checkliste S. 103: Tabellen müssen auf eine Seite passen, zentriert,
% umrandet und beschriftet (Überschrift oberhalb) sein.
% ============================================================

\chapter{Resultate}
\label{kap:resultate}

% Die Messzeitpunkt-Spalte steht in der CSV im Rohformat
% "YYYY-MM-DD HH:MM" und wird beim Kompilieren auf "DD.MM.YYYY HH:MM"
% umformatiert (Spaltenstil "ddmmyyyy from iso", definiert in
% preamble.tex).
Tabelle~\ref{tab:gaerung} zeigt den Verlauf der Hauptgärung anhand der
täglich protokollierten Messwerte.

\begin{table}[h]
  \centering
  \caption{Messprotokoll der Hauptgärung des Kanti Pale Ale.}
  \label{tab:gaerung}
  \footnotesize
  % Farben/Rahmen nach Vorlage tabelle-beispiel_2.html: weisse, dicke
  % Trennlinien zwischen allen Zellen (border-collapse-Optik), Kopfzeile
  % orange (bierheaderbg) mit schwarzer Schrift, Datenzeilen beige
  % (bierrowbg). Farben in preamble.tex definiert. extrarowheight sorgt
  % für mehr Abstand ober-/unterhalb des Zelltexts (auch bei mehrzeiligen
  % Bemerkung-Zellen, im Gegensatz zu \arraystretch).
  {\arrayrulecolor{white}\setlength{\arrayrulewidth}{1pt}%
  \setlength{\extrarowheight}{4pt}%
  \pgfplotstabletypeset[
    columns={Messzeitpunkt,Temperatur_Celsius,Spezifisches_Gewicht_SG,%
      Restzucker_Plato,Alkoholgehalt_Vol_Prozent,Bemerkung},
    columns/Messzeitpunkt/.style={
      ddmmyyyy from iso,
      column name={\bfseries Messzeitpunkt}, column type={|l|}},
    columns/Temperatur_Celsius/.style={
      column name={\bfseries Temp.\ (\si{\celsius})}, column type={c|}},
    columns/Spezifisches_Gewicht_SG/.style={
      column name={\bfseries Dichte (SG)}, column type={c|}},
    columns/Restzucker_Plato/.style={
      column name={\bfseries Restzucker (\si{\percent}P)}, column type={c|}},
    columns/Alkoholgehalt_Vol_Prozent/.style={
      column name={\bfseries Alkohol (\si{\percent}\,vol.)}, column type={c|}},
    columns/Bemerkung/.style={
      column name={\bfseries Bemerkung}, column type={p{3cm}|}},
    every head row/.style={before row={\hline\rowcolor{bierheaderbg}}, after row=\hline},
    every odd row/.style={before row=\rowcolor{bierrowbg}, after row=\hline},
    every even row/.style={before row=\rowcolor{bierrowbg}, after row=\hline},
    every last row/.style={after row=\hline},
  ]{daten/gaerung.csv}}%
  \par\vspace{0.3em}\footnotesize Eigene Messung, Kanti Pale Ale,
  August/Oktober 2026.
\end{table}

% Für das Diagramm werden die (bereits umformatierten) Messzeitpunkte
% zusätzlich als expliziete Tick-Label-Liste zusammengebaut: die
% Kurzform "xticklabels from table" übernimmt Zellinhalte 1:1 aus der
% Rohdatei und wendet "ddmmyyyy from iso" (ein reiner
% Typeset-/Einlese-Spaltenstil) dabei nicht an.
\pgfplotstableread[col sep=semicolon]{daten/gaerung.csv}\gaerungraw
\pgfplotstablegetrowsof{\gaerungraw}
\pgfmathtruncatemacro{\gaerungmaxrow}{\pgfplotsretval-1}
\foreach \gaerungrow in {0,...,\gaerungmaxrow}{%
  \pgfplotstablegetelem{\gaerungrow}{Messzeitpunkt}\of\gaerungraw
  \let\kswoRawDT\pgfplotsretval
  \StrLeft{\kswoRawDT}{4}[\kswoIsoY]
  \StrMid{\kswoRawDT}{6}{7}[\kswoIsoM]
  \StrMid{\kswoRawDT}{9}{10}[\kswoIsoD]
  \StrMid{\kswoRawDT}{12}{16}[\kswoIsoT]
  \edef\kswoNewDT{\kswoIsoD.\kswoIsoM.\kswoIsoY\space\kswoIsoT}
  \ifnum\gaerungrow=0
    \xdef\gaerungticklabels{\kswoNewDT}
  \else
    \xdef\gaerungticklabels{\gaerungticklabels,\kswoNewDT}
  \fi
}

Abbildung~\ref{fig:gaerverlauf} stellt denselben Datensatz grafisch dar
und zeigt, wie der Restzuckergehalt über den Gärverlauf sinkt, während
der Alkoholgehalt entsprechend ansteigt.

\begin{figure}[h]
  \centering
  \begin{tikzpicture}
    \begin{axis}[
      width=0.95\linewidth,
      height=6.5cm,
      xlabel={Messzeitpunkt},
      ylabel={Anteil (\si{\percent})},
      xtick=data,
      xticklabels/.expand once=\gaerungticklabels,
      x tick label style={rotate=60, anchor=east, font=\footnotesize},
      y tick label style={font=\footnotesize},
      label style={font=\footnotesize},
      legend style={font=\footnotesize},
      legend pos=north east,
      ymin=0, ymax=13,
      grid=major,
    ]
      \addplot table [x expr=\coordindex, y=Restzucker_Plato,
        col sep=semicolon] {daten/gaerung.csv};
      \addlegendentry{Restzucker (\si{\percent}P)}
      \addplot table [x expr=\coordindex, y=Alkoholgehalt_Vol_Prozent,
        col sep=semicolon] {daten/gaerung.csv};
      \addlegendentry{Alkohol (\si{\percent}\,vol.)}
    \end{axis}
  \end{tikzpicture}
  \caption{Verlauf von Restzucker und Alkoholgehalt während der
    Hauptgärung.}
  \label{fig:gaerverlauf}
\end{figure}


% ============================================================
% Diskussion (Handbuch S. 84)
%
% - Resultate mit bestehender Literatur/Erwartung vergleichen
% - Was ist neu? Welche Ergebnisse bedürfen weiterer Klärung?
% - Spekulationen sind am Schluss erlaubt, müssen aber klar als solche
%   gekennzeichnet sein
% - Fazit mit sachlichem Teil (zentrale Ergebnisse, offene Fragen) und
%   persönlichem Teil (eigene Einschätzung, was für künftige Projekte
%   gelernt wurde) -- je etwa gleich gewichtet (Handbuch S. 86)
%
% Beispiel Kurzbeleg gemäss gewähltem Zitierstil (Klammern), mit der
% echten Quelle aus quellen.bib:
%   Wörtliches Zitat: \parencite[S.~45]{bruecklmeier2022}
%   Paraphrase mit "Vgl.": \vgl[S.~45]{bruecklmeier2022}
% ============================================================

\chapter{Diskussion}
\label{kap:diskussion}



% ---- Quellenverzeichnis, aufgeteilt gemäss Handbuch Kap. 7.5.1 -----
% Jeder Eintrag in quellen.bib braucht das passende keywords-Feld
% (literatur / internet / abbildung / tabelle).
\printbibliography[keyword=literatur, title=Literatur]
\printbibliography[keyword=internet, title=Internetquellen]
\printbibliography[keyword=abbildung, title=Abbildungen]
\printbibliography[keyword=tabelle, title=Tabellen]

% ---- Verwendete KI-Tools, zwingend bei KI-Einsatz (Handbuch Kap. 7.4.2) ---
% ============================================================
% Verwendete KI-Tools (Handbuch Kap. 7.4.2, S. 92)
% "Die Verwendung von KI-Tools muss in einem zusätzlichen Abschnitt am
% Ende des normalen Quellenverzeichnisses nachgewiesen werden [...]:
% Name des Tools, Datum der Verwendung, eingegebene Aufforderungen
% (Prompts)." Zusätzlich ist ein Reflexionsformular einzureichen (siehe
% Hinweis in selbstaendigkeitserklaerung.tex).
%
% Diese Angaben nach Bedarf ergänzen, falls im Verlauf der Arbeit weitere
% KI-Tools oder Prompts hinzukommen.
% ============================================================

\chapter*{Verwendete KI-Tools}
\addcontentsline{toc}{chapter}{Verwendete KI-Tools}

Das vorliegende LaTeX-Dokumentgerüst (Formatierung gemäss Tab.\,14,
Kapitelstruktur gemäss Kap.\,7.1.2, Zitier- und
Bibliografie-Vorlagen gemäss Kap.\,7.5) wurde mit Unterstützung des
KI-Tools Claude erstellt. Der wissenschaftliche Inhalt der Arbeit
(Fragestellung, Durchführung, Auswertung, Diskussion) wird eigenständig
von den Autoren verfasst; der Beispieltext in der Einleitung wurde als
Platzhalter mit demselben Tool entworfen und ist vor Abgabe durch
eigene Formulierungen zu ersetzen.

\begin{itemize}
  \item \textbf{Name des Tools:} Claude Code (Claude Sonnet 5), Anthropic
  \item \textbf{Datum der Verwendung:} 10.\,September 2026
  \item \textbf{Eingegebene Aufforderungen (Auszug):}
    \begin{enumerate}
      \item[a)] \enquote{Baue mir, gemäss den Vorgaben im Handbuch, ein
        LaTeX-Grundgerüst für die Probearbeit zum Thema Bierbrauen
        (naturwissenschaftliche Arbeit, Kurzbelege in Klammern).}
      \item[b)] \enquote{Ergänze die Aufsätze/Zeitschriften-Treiber im
        Quellenverzeichnis nach demselben Muster.}
      \item[c)] \enquote{Schreibe einen kurzen Beispieltext in die
        Einleitung und füge eine Referenz auf das Buch
        \emph{bruecklmeier2022} ein.}
    \end{enumerate}
\end{itemize}


% ---- Anhang (optional) und Selbständigkeitserklärung (zwingend) ---
% ============================================================
% Anhang (optional, Handbuch S. 81)
%
% Nur Material aufnehmen, das im Haupttext mindestens einmal erwähnt
% wurde und das Verständnis der Arbeit unterstützt, ohne den Lesefluss
% zu unterbrechen. Für dieses Thema z.B.:
% - Detailliertes Brauprotokoll / Rezeptblatt
% - Rohdaten aus Messungen (z.B. Excel-Tabelle mit Gärverlauf)
% - Fotos vom Brauprozess
% - Fragebogen samt Rohdaten, falls eine Degustationsumfrage gemacht wurde
% ============================================================

\appendix
\chapter{Anhang}
\label{kap:anhang}


% ============================================================
% Selbständigkeitserklärung (zwingend, Handbuch S. 94)
% Wortlaut wörtlich aus dem Handbuch übernommen. Da die Probearbeit zu
% zweit verfasst wird, erklärt und unterschreibt jede Autorin / jeder
% Autor die Erklärung einzeln (die Vorlage verwendet "ich", nicht "wir").
%
% Muss handschriftlich unterschrieben werden -- nach dem Ausdrucken von
% Hand ausfüllen.
% ============================================================

\chapter*{Selbständigkeitserklärung}
\addcontentsline{toc}{chapter}{Selbständigkeitserklärung}

\noindent
Ich erkläre hiermit, dass ich
\begin{itemize}
  \item diese Arbeit ohne unerlaubte fremde Hilfe angefertigt habe.
  \item keine anderen als die angegebenen Quellen benutzt habe.
  \item sämtliche Stellen, die wörtlich aus fremden Quellen entnommen
    oder mit eigenen Worten wiedergegeben wurden, als solche
    gekennzeichnet und die Urheberschaft angegeben habe.
  \item die Meinungen und Ideen anderer explizit ausgewiesen habe.
  \item die Arbeit in eigenen Worten formuliert habe.
\end{itemize}

\noindent
Ich bin damit einverstanden, dass meine Arbeit als pdf-Datei in der
Mediothek archiviert wird und Schulangehörige diese einsehen dürfen.

% Bei Verwendung von KI-Tools (siehe ki-tools.tex) muss gemäss Handbuch
% Kap. 7.4.2 zusätzlich ein Reflexionsformular eingereicht werden --
% siehe reflexionsformular.tex (eigenständiges Dokument, jede/r Autor/in
% füllt ein eigenes Exemplar aus).

\vspace{3cm}

\noindent
\begin{minipage}{0.45\textwidth}
  Ort, Datum \\[2.5cm]
  \hrulefill \\
  Vorname Nachname 1
\end{minipage}
\hfill
\begin{minipage}{0.45\textwidth}
  Ort, Datum \\[2.5cm]
  \hrulefill \\
  Vorname Nachname 2
\end{minipage}


\end{document}
